\begin{figure}[!t]
\centering
\begin{tikzpicture}
\begin{axis}[axstd, height=6.0cm,
    xlabel={Offered load (\% of service capacity)},
    ylabel={$\Prb(L > \tau_{\max})$},
    xmin=5, xmax=95,
    ymode=log, ymin=2e-5, ymax=2,
    xtick={10,20,30,40,50,60,70,80,90},
    legend pos=north west]
\addplot[b1] coordinates {\Esixstatic}; \addlegendentry{Always-purify}
\addplot[sp] coordinates {\Esixfull};   \addlegendentry{\SP{} (adaptive gate $q=0.31$)}
\addplot[b2] coordinates {\Esixundef};  \addlegendentry{Undefended}
\end{axis}
\end{tikzpicture}
\caption{Deadline-violation probability from a discrete-event
queueing simulation driven by measured service times
($1.21$\,ms decode, $8.06$\,ms purification at $K=4$,
$\tau_{\max}=25$\,ms). Below $70\%$ load the undefended arm is
strictly better---security costs tail latency. Above $80\%$ it
collapses, because corrupted frames provoke retransmissions that
compound the congestion causing them.}
\label{fig:e6}
\end{figure}
